% ============================================================
% AUTOMATISCH GEGENEREERD BESTAND
% Niet handmatig aanpassen.
% ============================================================

\ifdefined\HCode
  % Geen afzonderlijke module-openingspagina in HTML.
\else
  \FVDmoduleopening
    {Van data naar functies — verbanden herkennen en gedrag onderzoeken}
    {In deze module leggen we de basis voor het onderzoeken van verbanden en functies. We starten met een korte herhaling van belangrijke reken- en algebraïsche vaardigheden die we doorheen de verdere cursus nodig hebben. Daarna onderzoeken we verbanden tussen grootheden aan de hand van gegevens. Met spreidingsdiagrammen, correlatie en lineaire regressie leren we patronen herkennen en modellen beoordelen. Daarbij besteden we bijzondere aandacht aan het verschil tussen samenhang en causaliteit. Vervolgens maken we de overstap van relaties naar functies. We bouwen de taal op waarmee we functies kunnen beschrijven en onderzoeken: domein en bereik, nulwaarden, tekenverloop, stijgen en dalen, extrema, symmetrie en verschillende voorstellingsvormen. Grafieken gebruiken we niet alleen om functies te visualiseren, maar ook om vergelijkingen en ongelijkheden op te lossen en om het gedrag van functies te onderzoeken. Ten slotte bestuderen we wat er gebeurt wanneer een veranderlijke een punt nadert of onbegrensd groot wordt. Zo ontstaan de begrippen limiet, continuïteit en asymptotisch gedrag. Verbanden en functies vormen een centrale taal van de wiskunde. Dezelfde ideeën gebruiken we in onder andere fysica, chemie, biologie, data-analyse, engineering en andere wetenschappen om verschijnselen te beschrijven, modellen op te stellen en voorspellingen te maken.}
    {%
    \FVDmodulethemetile{1}{Parate kennis — Het gereedschap dat je nodig hebt}
    \FVDmodulethemetile{2}{Van werkelijkheid naar wiskunde — en terug}
    \FVDmodulethemetile{3}{Verbanden onderzoeken met data}
    \FVDmodulethemetile{4}{Van relatie naar functie}
    \FVDmodulethemetile{5}{Het gedrag van functies beschrijven}
    \FVDmodulethemetile{6}{Functies in verschillende vormen}
    \FVDmodulethemetile{7}{Vergelijkingen door een functiebril}
    \FVDmodulethemetile{8}{Wat gebeurt er dichtbij een punt?}
    \FVDmodulethemetile{9}{Gedrag op oneindig en asymptoten}
    \FVDmodulethemetile{10}{Continuïteit en limieten}
    \FVDmodulethemetile{11}{Functies analyseren in context}
    }
\fi
